Este reporte estudia el problema de encontrar el camino de mayor puntaje para
un rover en una cuadrícula de $N\times M$ celdas con murallas y minas, donde
solo son posibles los movimientos hacia la derecha y hacia abajo y el rover
dispone de $V$ vidas. Se comparan tres paradigmas: fuerza bruta, programación
dinámica (DP) y un algoritmo greedy. El objetivo es relacionar el tiempo y la
memoria medidos con la complejidad teórica de cada uno, y evaluar la calidad
de las soluciones de greedy respecto al óptimo mediante la brecha
$\Delta = s(\mathrm{OPT}) - s(\mathrm{Greedy})$ y su tasa de falla.